%!TEX root = ../chapter05.tex 
%----------------------------------------------------------------------------------------
%	
%----------------------------------------------------------------------------------------
%----------------------------------------------------------------------------------------
%	 
%----------------------------------------------------------------------------------------
\newpage 
\loadgeometry{marginlayout}  
\pagestyle{scrheadings} % Print headers again  
\KOMAoptions{
	headwidth=textwithmarginpar,
	headsepline=true,
	headsepline= 0.5pt,
	footwidth=textwithmarginpar,
}   
\onehalfspacing  
\section{Triangles rectangles}
\begin{definition}
	L'hypoténuse d'un triangle rectangle est le côté opposé à  l'angle droit.
\end{definition}
\begin{marginfigure}[-1cm]
	\begin{tikzpicture}[scale=1.0]
		\tkzDefPoint(0,0){A} \tkzDefPoint(4,0){B}
		\tkzDefTriangle[golden](A,B)\tkzGetPoint{C}
		\tkzDrawPolygon(A,B,C) \tkzDrawPoints(A,B,C)
		\tkzLabelPoints[below left](A)
		\tkzLabelPoints[below right](B)
		\tkzLabelPoints[above](C)
		%		\tkzMarkAngle[ arc=l, size=0.5](B,A,C)
		\tkzMarkRightAngle[ pos=.15, size=0.4](C,B,A)
		%		\tkzLabelAngle[pos=1.5,circle](B,A,C){$\alpha$}
		%		\tkzLabelSegment[below](A,B){\small adjacent à $\alpha$}
		\tkzLabelSegment[ above , sloped](A,C){hypoténuse}
		%		\tkzLabelSegment[below, sloped](B,C){opposé à $\alpha$}
	\end{tikzpicture} 
\end{marginfigure}
\begin{theorem}[Théorème de Pythagore]  
	Dans un triangle rectangle, le carré de l'hypoténuse est égal à  la somme des carrés des côtés de l'angle droit.
\end{theorem}
\begin{remark} Conséquence : l'hypoténuse est bien le plus grand côté.
	\textcolor{white}{
		\begin{align*}
			AB^2 + BC^2 & = AC^2 && \text{$AB$, $BC$ et $AC\geqslant0$}\\
			AB^2 & \leqslant AC^2 \\
			AB & \leqslant AC
		\end{align*}
	}
\end{remark}

%----------------------------------------------------------------------------------------
%	 
%----------------------------------------------------------------------------------------
%\newpage 
\begin{marginfigure} 
	\begin{tikzpicture}[scale=1]
		\tkzSetUpPoint[size=3, fill=black]
		\tkzfctset{tan style/.style={->,>=latex, color=red, line width=1.2pt}}   
		\tkzSetUpLine[line width= 1.2pt, add=.5 and .5] 
		%		 
		\tkzDefPoint(0,0){B}  
		\tkzDefShiftPoint[B](40:3.0){A}
		\tkzDefShiftPoint[B](15:4.0){C}
		\tkzDrawLine[add=0.2 and 0.2](B,C)
		\tkzDefLine[orthogonal=through A](B,C) \tkzGetPoint{a}
		\tkzLabelLine[pos=1.25,above left](B,C){$(\Delta)$}
		\tkzInterLL(B,C)(A,a) \tkzGetPoint{H}
		\tkzDrawLine[dashed](A,H)
		\tkzMarkRightAngles[size=0.2](C,H,A)
		\tkzDrawPoints(A,H)
		\tkzLabelPoints[below left,font=\scriptsize](A)  
		\tkzLabelPoints[below right,font=\scriptsize](H)   
	\end{tikzpicture} 
\end{marginfigure}  
\begin{definition} Soit une droite $(\Delta)$ et un point $A$ du plan. Le \textbf{projeté orthogonal} $H$ de $A$ sur $(\Delta)$ est le point d'intersection de $\Delta$ et de la perpendiculaire à $(\Delta)$ passant par $A$.
\end{definition} 
\begin{theorem}
	Soit une droite $(\Delta)$ et $A$ un point n'appartenant pas à $(\Delta)$. $H$ le projeté orthogonal de $A$ sur $(\Delta)$.\\
	La distance ente $A$ et la droite $(\Delta)$ est égale à $AH$. C'est la plus petite distance entre $A$ et un point de la droite $(\Delta)$. 
\end{theorem} 
\begin{proof} \emph{au programme} ~ Pour tout point $M\in(Delta)$, $AMH$ est un triangle rectangle d'hypoténuse $AM$, et $AM\geqslant AH$.\hfill 
	\begin{marginfigure}[-1cm] 
		\captionsetup{type=figure} 
		\begin{tikzpicture}[scale=1]
			\tkzSetUpPoint[size=3, fill=black]
			\tkzfctset{tan style/.style={->,>=latex, color=red, line width=1.2pt}}   
			\tkzSetUpLine[line width= 1.2pt, add=.5 and .5] 
			%		 
			\tkzDefPoint(0,0){M}  
			\tkzDefShiftPoint[M](45:3.0){A}
			\tkzDefShiftPoint[M](15:4.0){C}
			\tkzDrawLine[add=0.2 and 0.2](B,C)
			\tkzDefLine[orthogonal=through A](M,C) \tkzGetPoint{a}
			\tkzLabelLine[pos=1.25,above left](M,C){$(\Delta)$}
			\tkzInterLL(M,C)(A,a) \tkzGetPoint{H}
			\tkzDrawLine[dashed](A,H)
			\tkzMarkRightAngles[size=0.2](C,H,A)
			\tkzDrawPoints(A,H,M)
			\tkzLabelPoints[above right,font=\scriptsize](A)  
			\tkzLabelPoints[below right,font=\scriptsize](H)   
			\tkzLabelPoints[below,font=\scriptsize](M)  
			\tkzDrawSegments[dashed](A,M)
			\tkzLabelSegment[font=\scriptsize](A,H){$\delta$}
		\end{tikzpicture} 
	\end{marginfigure}
	
	
%	\vfill 
	
\end{proof}
% \begin{remark} On note que le cercle de centre $A$ et de rayon $AH$ touche la droite $(\Delta)$ en un unique point $H$. On dit que la droite $(\Delta)$ est tangente au cercle.
% \end{remark}
% \begin{definition}	Soit un cercle de centre $O$ et de rayon $R$. En chacun de ses points, le cercle admet une tangente. 
% 	
% 	La tangente en $M$ au cercle est la droite passant par $M$ et perpendiculaire au rayon $OM$.
% 	
% 	La distance entre une tangente et le centre du cercle est égale au rayon $R$.
% \end{definition}
\begin{figure}[hb]\centering 
	\begin{tikzpicture}[scale=1]
		\tkzSetUpPoint[size=3, fill=black]
		\tkzfctset{tan style/.style={->,>=latex, color=red, line width=1.2pt}}   
		\tkzSetUpLine[line width= 1.2pt, add=.5 and .5] 
		%		 
		\tkzDefPoint(0,0){O}
		\tkzDefPoint(3,3){E}
		\tkzDefShiftPoint[O](13:2.5){M}
		%		\tkzDefRandPointOn[circle=center O radius 2.5cm]
		%		\tkzGetPoint{M}
		\tkzDefShiftPoint[O](150:2.5){N}
		%		\tkzDefRandPointOn[circle=center O radius 2.5cm]
		%		\tkzGetPoint{N}
		
		\tkzDefShiftPoint[O](-50:2.5){A}
		\tkzDefShiftPoint[O](-120:2.5){B}
		\tkzDrawSegments(O,M O,N)
		\tkzDrawCircle(O,M)
		\tkzDefTangent[at=M](O)
		\tkzGetPoint{h}
		\tkzDefTangent[at=N](O)
		\tkzGetPoint{i}
		\tkzDrawLine[add = 2 and 1.5](M,h) 
		\tkzLabelLine[pos=1.25,above left](M,h){$(\Delta)$}
		\tkzDrawLine[add = 2 and 2](N,i)
		\tkzLabelLine[pos=1.25,above left](N,i){$(\Delta')$}
		\tkzDrawLine[dashed, add = 0.5 and 0.5](A,B)
		\tkzMarkRightAngle[](O,M,h)
		\tkzMarkRightAngle[](O,N,i)
		\tkzLabelPoints[above right,font=\scriptsize](O)  
		\tkzLabelPoints[below right,font=\scriptsize](M)  
		\tkzLabelPoints[above left,font=\scriptsize](N)  
		\tkzLabelPoints[below left,font=\scriptsize](B)  
		\tkzLabelPoints[below right,font=\scriptsize](A)  
		\tkzLabelSegments[font=\scriptsize](O,N O,M){$R$}    
		\tkzLabelCircle[font=\scriptsize,R, above](O,2.5)(90){$\mathscr{C}$}
	\end{tikzpicture} 
	\caption{La droite $(AB)$ est sécante au cercle $\mathscr{C}$.\\
		Les droites $(\Delta)$ et $(\Delta')$ sont dites tangentes au cercle. \\
		La tangente en $M$ au cercle en $M$ est perpendiculaire au rayon $OM$.\\
		La distance entre une tangente et le centre du cercle est égale au rayon $R$.\\
		$M$ est le projeté orthogonal du centre $O$ sur droite $(\Delta)$.
	} 
\end{figure}

%----------------------------------------------------------------------------------------
%	
%----------------------------------------------------------------------------------------
\newpage 



\begin{definition} Dans un triangle rectangle d'hypoténuse $1$. $\alpha$ un angle aigu.  On designe par :
	\begin{enumerate}[label=\textbullet, leftmargin=*]
		\item $\cos(\alpha)$ la longueur du côté adjacent à $\alpha$
		\item $\sin(\alpha)$ la longueur du côté opposé à $\alpha$
	\end{enumerate}
	
	\begin{marginfigure} 
	\begin{tikzpicture}[scale=1.0]
		\tkzDefPoint(0,0){A} \tkzDefPoint(4,0){B}
		\tkzDefTriangle[golden](A,B)\tkzGetPoint{C}
		\tkzDrawPolygon(A,B,C) \tkzDrawPoints(A,B,C)
		\tkzLabelPoints[below left](A)
		\tkzLabelPoints[below right](B)
		\tkzLabelPoints[above](C)
		\tkzMarkAngle[mark=none, arc=l, size=0.5](B,A,C)
		\tkzMarkRightAngle[ pos=.0.15, size=0.4](C,B,A)
		\tkzLabelAngle[pos=1.5,circle](B,A,C){$\alpha$}
		\tkzLabelSegment[below](A,B){\small  $\cos(\alpha)$}
		\tkzLabelSegment[ above , sloped](A,C){$1$}
		\tkzLabelSegment[below, sloped](B,C){\small  $\sin(\alpha)$}
	\end{tikzpicture}  
\end{marginfigure} 
\end{definition}    
\begin{theorem} Pour tout valeur de l'angle $\alpha$ on a :
\[
 \cos^2(\alpha) + \sin^2(\alpha) = 1
\] 
\end{theorem}
\begin{proof}[démonstration] (\emph{au programme}) 
	Conséquence directe du théorème de Pythagore ! 
\end{proof}
\begin{marginfigure}[2cm]
	\begin{tikzpicture}[scale=1.0]
		\tkzDefPoint(0,0){A} \tkzDefPoint(4,0){B}
		\tkzDefTriangle[golden](A,B)\tkzGetPoint{C}
		\tkzDrawPolygon(A,B,C) \tkzDrawPoints(A,B,C)
		\tkzLabelPoints[below left](A)
		\tkzLabelPoints[below right](B)
		\tkzLabelPoints[above](C)
		\tkzMarkAngle[ mark=none,arc=l, size=0.5](B,A,C)
		\tkzMarkRightAngle[ pos=.15, size=0.4](C,B,A)
		\tkzLabelAngle[pos=1.5,circle](B,A,C){$\alpha$}
		\tkzLabelSegment[below](A,B){\small adjacent à  $\alpha$}
		\tkzLabelSegment[ above , sloped](A,C){hypoténuse}
		\tkzLabelSegment[below, sloped](B,C){opposé à  $\alpha$}
	\end{tikzpicture}  
	\caption{En 2\up{nd}, on calcule des rapports trigonométriques d'angles aigus.}
\end{marginfigure}
\begin{definition}
	Soit le triangle $ABC$ rectangle en $B$. 
	\begin{description} 
		\item[le \textbf{sinus}]   de l'angle $\alpha$   
		$$ \sin(\alpha) = \frac{\text{côté opposé à  $\alpha$}}{\text{hypoténuse}} = \frac{BC}{AC}  \leqslant  1$$ 
		\item[le \textbf{cosinus}] de l'angle $\alpha$ : 
		$$ \cos(\alpha) = \frac{\text{côté adjacent à  $\alpha$}}{\text{hypoténuse}} = \frac{AB}{AC} \leqslant  1 $$ 
		\item[la \textbf{tangente}] de l'angle $\alpha$ :
		$$ \tan(\alpha) = \frac{\text{côté opposé à  $\widehat{BAC}$}}{\text{côté adjacent à  $\alpha$}} = \frac{BC}{AB} $$
	\end{description} 
\end{definition}
 
		\vfill 